\section{Introduction}\label{sec:introduction}

Global solvers for nonconvex mixed-integer nonlinear programs (MINLPs) bound
the optimal value by relaxing every nonlinear expression. The standard
construction decomposes each expression into elementary operations,
introduces an auxiliary variable for each, and relaxes each operation on its
own \citep{McCormick1976,SmithPantelides1999,TawarmalaniSahinidis2004}. This
construction is general, but it relaxes expressions as if they were
unrelated. In many models they are not: the same variables appear in several
nonlinear constraints, and linear constraints restrict how these variables
can move together. Pooling and blending models, heat and mass exchange
networks and gas networks are typical examples. A relaxation that treats such
expressions separately can be much weaker than one that treats them jointly.

The strongest convex relaxation of several functions $f_1,\dots,f_p$ of the
same variables $x$ on a common domain $D$ is the convex hull of their joint
graph $\{(x,f_1(x),\dots,f_p(x)):x\in D\}$. Its valid linear inequalities are
\emph{support cuts}: for a direction $(a,b)$, the inequality
$a\T x+b\T w\ge\min_{x\in D}\{a\T x+b\T F(x)\}$ is valid, where $w_j$ stands for
$f_j(x)$, and all such inequalities together describe the hull. This
simultaneous convexification has been studied for vectors of univariate
functions, for multilinear, bilinear and composite functions, and in
applications \citep{Tawarmalani2010,Ballerstein2013,LiersEtAl2021,DavarniaRichardTawarmalani2017,HeTawarmalani2021,ZhuHeTawarmalani2026}.
Turning it into a solver component raises four questions.
\begin{enumerate}[label=(\arabic*)]
\item \emph{Which blocks are worth relaxing jointly?} In particular, is it
enough to relax every pair of interacting variables exactly, or do larger
blocks carry information that pairs miss?
\item \emph{Exactness.} The support value is the minimum of a nonconvex
function. For which blocks can it be computed exactly, so that the cuts are
as strong as possible?
\item \emph{Integration and validity.} Support cuts live in the space of
$(x,w)$, and introducing the variables $w$ changes the model that the solver
presolves, propagates and branches on. Can the cuts be stated in the original
variables? And since directions are found numerically and the solver stores
binary64 rows, which part of the computation must be exact for the stored row
to be valid?
\item \emph{Usefulness.} Modern solvers already detect quadratic, bilinear and
convex structure and have their own cut families for it
\citep{BestuzhevaEtAl2025}. Do joint support cuts improve them?
\end{enumerate}

We answer these questions for quadratic blocks in depth and for shared
nonlinear expressions more generally. The contributions are the following.

\begin{itemize}
\item \emph{Limits of composing pair hulls} (Section~\ref{sec:composition}). On
a path of three variables, exact pair hulls identified on the first two
moments of the shared variable miss a gap of $1/128$ that a single support
cut in the existing coordinates closes. For a family of such paths the gap is
positive exactly when two two-point sets interleave, and it then equals
$\delta^2/2$, where $\delta$ is their distance. An interface that identifies
$\kappa$ moments of the shared variable gives the trivial bound zero exactly
when the local zero sets alternate at least $\kappa+1$ times, a reading of the
classical description of Radon partitions on the moment curve, and for every
$\kappa$ there is a quadratic instance with a relative gap of $100\%$. A dense
semidefinite relaxation with the McCormick inequalities of the missing
product closes every three-variable path, by a theorem of
\citet{BurerNatarajanWillemsen2025}; the advantage of joint blocks over their
pairs is therefore one of representation, which support cuts obtain in the
coordinates that the model already has. We also relate the gain from merging
pairs into a star to the glued relaxation.
\item \emph{Exact support for constrained stars} (Section~\ref{sec:quadratic}).
For quadratic stars whose affine rows couple the center with one leaf at a
time, we give an exact rational algorithm with $O(m_0+(m+k)\log(m+k))$
operations for $k$ leaves, $m$ center--leaf rows and $m_0$ rows in the center
alone; on instances with $m=\Theta(k)$ this is optimal for algebraic
computation trees. The box-constrained case is a special case of the forest
algorithm of \citet{DelPiaKhajavirad2026}; to our knowledge, no exact
algorithm for the case with center--leaf rows has been published. For
polytopes of fixed dimension we state face enumeration together with the
degenerate cases that a certificate must cover.
\item \emph{A certified separator and its evaluation}
(Sections~\ref{sec:original}--\ref{sec:computations}). Minimizing nonnegative
combinations of nonlinear rows over their common block domain and eliminating
the rows gives cuts in the original variables, with no auxiliary variables.
Together these cuts describe the projection of the joint-graph relaxation,
and we give conditions under which this projection is the convex hull of the
set that the rows define, and examples where it is not. Validity of a stored
cut rests on four checks: a certified lower bound over the whole block domain
for the exact binary64 direction, exact elimination, bound-corrected rounding,
and a comparison with the row that the solver stores; the direction itself may
be found by any numerical method. We implemented the cuts as a SCIP separator
whose model import preserves the exact binary64 semantics and the domains of
the source model, and whose every cut can be replayed. We evaluate it with
prospective protocols against SCIP's default, against SCIP's
disabled-by-default nonconvex separators, and, on a constructed family,
against Gurobi.
\end{itemize}

The other ingredients, namely the support description, Bernstein and
ball-arithmetic bounds, safe rounding, Lagrangian duality and oracle
separation, are classical; we use them and state them precisely where a
certificate depends on the details (Sections~\ref{sec:original}
and~\ref{sec:certification}, Appendices~\ref{app:bounds}
and~\ref{app:separation}).

Section~\ref{sec:computations} reports four campaigns, each run with code,
limits and protocol fixed in advance. All 143{,}267 recorded cuts of the last
two campaigns passed a fresh-process replay, which re-executes the exact
support routines against the source model and the row stored by SCIP.
Constants computed without a certificate, from samples or from a local solver,
would have been wrong for 38\% of the cuts on MINLPLib models and for nearly
all cuts on the path family, and would have cut off feasible solutions. On MINLPLib models, including 20 larger models that SCIP
does not solve within 60 seconds, the cuts did not change which models SCIP
solved in campaigns 2 to~4 and left SCIP's own solving time unchanged; the
Python separator made the runs slower, and SCIP's own disabled-by-default
separators improved the root bound more often. On constructed instances with
the path structure of Section~\ref{sec:composition}, SCIP's root bound is weak
because of a presolve step; with that step disabled, SCIP reaches the level of
glued pair relaxations and no further. A post hoc diagnostic, confirmed by a
later prospective control, showed that the separator closes the gap beyond
that level only if it may add enough cuts per block and tries the exact
support of each whole row first. On 40 fresh
instances, half of them with a binding coupling row, this configuration, fixed
before the instances were generated, closed at least 97\% of SCIP's root gap
on every instance and solved all 40 within 300 seconds, against 19 for SCIP,
22 for SCIP with its disabled separators and 11 for Gurobi.


We do not claim a new convexification principle, a general speedup, or a
certificate for a complete solve; the certificates cover the cuts and the
model they were derived from. Section~\ref{sec:setting} fixes notation and
reviews related work, Sections~\ref{sec:composition}--\ref{sec:certification}
develop the theory, Section~\ref{sec:implementation} describes the
implementation, Section~\ref{sec:computations} reports the experiments, and
Section~\ref{sec:conclusions} discusses consequences and open problems. The
appendices contain proofs, the first two campaigns, complete separation, and
the test sets.
